% answers_short.tex —— 2016 年·天津卷·理科 · 简短答案速查
% 由 tools/make_short_answers.py 自 answers.tex 的 \gksolhead/\gkanswer 提取，
% 并与 2016/2016_数学真题_天津理_答案要点.md 逐题比对后生成，勿手改（改源 MD 后重跑）。

\begin{gknotice}
\textbf{使用说明}：本节只给最终答案便于快速核对，完整推导见解析版。答案为编写者独立推导，\textbf{非官方标准答案}。
\end{gknotice}

\gkpart{选择题}
\begin{center}\begin{tabular}{ccccccccc}
\toprule
题号 & 1 & 2 & 3 & 4 & 5 & 6 & 7 & 8 \\
\midrule
答案 & D & B & A & B & C & D & B & C \\
\bottomrule
\end{tabular}\end{center}

\gkpart{填空题}
\begin{itemize}[leftmargin=2.6em]
  \item[\textbf{9.}] \(2\)
  \item[\textbf{10.}] \(-56\)
  \item[\textbf{11.}] \(2\)（\(\mathrm{m}^{3}\)）
  \item[\textbf{12.}] \(\dfrac{2\sqrt{3}}{3}\)
  \item[\textbf{13.}] \(\left(\dfrac{1}{2},\dfrac{3}{2}\right)\)
  \item[\textbf{14.}] \(\sqrt{6}\)
\end{itemize}

\gkpart{解答题}
\begin{itemize}[leftmargin=2.6em]
  \item[\textbf{15.}] (1) 定义域 \(\left\{x\mid x\neq\dfrac{\pi}{2}+k\pi,\ k\in\mathbf{Z}\right\}\)，最小正周期 \(T=\pi\)；(2) 在 \(\left[-\dfrac{\pi}{4},-\dfrac{\pi}{12}\right]\) 上单调递减，在 \(\left[-\dfrac{\pi}{12},\dfrac{\pi}{4}\right]\) 上单调递增
  \item[\textbf{16.}] (1) \(P(A)=\dfrac{1}{3}\)；(2) \(P(X=0)=\dfrac{4}{15}\)，\(P(X=1)=\dfrac{7}{15}\)，\(P(X=2)=\dfrac{4}{15}\)；\(E(X)=1\)
  \item[\textbf{17.}] (1) 见解析版；(2) \(\dfrac{\sqrt{3}}{3}\)；(3) \(\dfrac{\sqrt{7}}{21}\)
  \item[\textbf{18.}] (1) 见解析版；(2) 见解析版
  \item[\textbf{19.}] (1) \(\dfrac{x^{2}}{4}+\dfrac{y^{2}}{3}=1\)；(2) \(k\in\left(-\infty,-\dfrac{\sqrt{6}}{4}\right]\cup\left[\dfrac{\sqrt{6}}{4},+\infty\right)\)
  \item[\textbf{20.}] (1) \(a\leqslant0\) 时在 \(\mathbf{R}\) 上单调递增；\(a>0\) 时在 \((-\infty,1-\sqrt{\frac{a}{3}}]\)、\([1+\sqrt{\frac{a}{3}},+\infty)\) 上递增，在 \([1-\sqrt{\frac{a}{3}},1+\sqrt{\frac{a}{3}}]\) 上递减；(2) 见解析版；(3) 见解析版
\end{itemize}
